\documentclass[12pt]{article}
\usepackage{amsmath, amssymb}
\usepackage[margin=1in]{geometry}
\setlength{\parskip}{6pt}
\title{QUBO Formulation: Subset Sum Problem}
\date{}
\begin{document}
\maketitle

\subsection*{Subset Sum Problem}

\paragraph{Problem Statement.}
Given a set of $N$ positive integers $\{a_1, a_2, \dots, a_N\}$ and a target value $T$, the Subset Sum Problem seeks to identify a subset of these integers whose sum is as close as possible to $T$. Formally, the goal is to determine a selection of indices such that the absolute difference between the sum of the selected integers and the target value is minimized.

\paragraph{Decision Variables.}
We introduce binary decision variables $x_i \in \{0, 1\}$ for each $i \in \{1, \dots, N\}$. The variable $x_i$ indicates whether the $i$-th integer $a_i$ is included in the subset:
\begin{equation}
    x_i = \begin{cases}
    1 & \text{if } a_i \text{ is selected}, \\
    0 & \text{otherwise}.
    \end{cases}
\end{equation}
The sum of the selected integers is represented by $\sum_{i=1}^N a_i x_i$.

\paragraph{QUBO Derivation.}
The objective is to minimize the absolute deviation between the subset sum and the target. This can be expressed as:
\begin{equation}
    \min \left| \sum_{i=1}^N a_i x_i - T \right|.
\end{equation}
Since the square function is monotonically increasing for non-negative arguments, minimizing the absolute deviation is equivalent to minimizing the squared deviation:
\begin{equation}
    \min \left( \sum_{i=1}^N a_i x_i - T \right)^2.
\end{equation}
We expand the squared term algebraically:
\begin{align}
    \left( \sum_{i=1}^N a_i x_i - T \right)^2 &= \left( \sum_{i=1}^N a_i x_i \right)^2 - 2T \sum_{i=1}^N a_i x_i + T^2 \\
    &= \sum_{i=1}^N \sum_{j=1}^N a_i a_j x_i x_j - 2T \sum_{i=1}^N a_i x_i + T^2.
\end{align}
Separating the terms where $i = j$ from those where $i \neq j$, and utilizing the idempotent property of binary variables $x_i^2 = x_i$, we obtain:
\begin{align}
    \sum_{i=1}^N \sum_{j=1}^N a_i a_j x_i x_j &= \sum_{i=1}^N a_i^2 x_i^2 + \sum_{i \neq j} a_i a_j x_i x_j \\
    &= \sum_{i=1}^N a_i^2 x_i + \sum_{i \neq j} a_i a_j x_i x_j.
\end{align}
The double sum over $i \neq j$ can be rewritten as twice the sum over ordered pairs with $i < j$:
\begin{equation}
    \sum_{i \neq j} a_i a_j x_i x_j = 2 \sum_{1 \le i < j \le N} a_i a_j x_i x_j.
\end{equation}
Substituting this back into the expansion yields the Hamiltonian $H(x)$:
\begin{equation}
    H(x) = \sum_{i=1}^N (a_i^2 - 2T a_i) x_i + 2 \sum_{1 \le i < j \le N} a_i a_j x_i x_j + T^2.
\end{equation}
From this general form, we read off the QUBO matrix entries $Q_{i,j}$ and the constant offset $c$. The QUBO objective is to minimize $x^\top Q x + c$, where $x$ is the vector of variables. The entries are given by:
\begin{align}
    Q_{i,i} &= a_i^2 - 2T a_i, \quad \forall i \in \{1, \dots, N\}, \\
    Q_{i,j} &= 2 a_i a_j, \quad \forall 1 \le i < j \le N, \\
    c &= T^2.
\end{align}
Note that $Q_{i,j} = Q_{j,i}$ for $i > j$ due to symmetry, and the off-diagonal terms for $i \neq j$ are captured by the factor of 2 in the summation.

\paragraph{Penalty Weight Justification.}
The Subset Sum Problem formulation above contains no constraints; any assignment of binary values to $x_i$ corresponds to a valid subset. Consequently, no penalty terms are required, and no penalty weights $\lambda$ are needed. The Hamiltonian consists solely of the objective function terms.

\paragraph{Correctness Argument.}
Minimizing $H(x)$ over $x \in \{0, 1\}^N$ recovers the optimal solution because $H(x)$ is exactly the squared deviation $(\sum a_i x_i - T)^2$ for any binary vector $x$. Since the squared deviation is non-negative and zero if and only if the subset sum equals the target, the global minimum of $H(x)$ corresponds to the subset with the minimum squared deviation. Because the square function preserves the ordering of non-negative values, the subset minimizing the squared deviation also minimizes the absolute deviation. Thus, the configuration $x^*$ that minimizes $H(x)$ identifies the subset of integers whose sum is closest to $T$.

\paragraph{Illustrative Example.}
Consider an instance with $N=3$ integers, values $a_1 = 1$, $a_2 = 2$, $a_3 = 3$, and target $T = 4$.
The general formulas yield the following QUBO parameters:
\begin{align}
    Q_{1,1} &= 1^2 - 2(4)(1) = -7, \\
    Q_{2,2} &= 2^2 - 2(4)(2) = -12, \\
    Q_{3,3} &= 3^2 - 2(4)(3) = -15, \\
    Q_{1,2} &= 2(1)(2) = 4, \\
    Q_{1,3} &= 2(1)(3) = 6, \\
    Q_{2,3} &= 2(2)(3) = 12, \\
    c &= 4^2 = 16.
\end{align}
The resulting Hamiltonian is:
\begin{equation}
    H(x) = -7x_1 - 12x_2 - 15x_3 + 4x_1x_2 + 6x_1x_3 + 12x_2x_3 + 16.
\end{equation}
Evaluating $H(x)$ for all $2^3 = 8$ configurations reveals the minimum energy. The configuration $x^* = (1, 0, 1)$ yields:
\begin{equation}
    H(1, 0, 1) = -7(1) - 12(0) - 15(1) + 4(0) + 6(1) + 12(0) + 16 = -22 + 6 + 16 = 0.
\end{equation}
This corresponds to selecting $a_1$ and $a_3$, giving a sum of $1 + 3 = 4$, which matches the target exactly with zero deviation. The minimum energy is $0$, confirming the optimal subset.
\end{document}